\chapter{Conclusão}

Este trabalho teve como objetivo analisar o consumo de energia da Java Virtual
Machine (JVM) por meio de um conjunto diversificado de \textit{benchmarks},
considerando separadamente energia e tempo de execução. Para isso, foi
desenvolvida uma extensão ao DaCapo Chopin \cite{dacapochopin} capaz de
coletar medições energéticas com o apoio do JRAPL \cite{jrapl}, o que tornou
possível integrar à execução dos \textit{benchmarks} métricas que não
estão presentes originalmente. Essa extensão constitui uma das
principais contribuições do trabalho, pois amplia o uso do DaCapo como
ferramenta para estudos de eficiência energética na JVM.

Os resultados obtidos mostram que consumo de energia e tempo de execução, ainda
que relacionados, não se comportam de forma equivalente. De modo geral, as
configurações de menor tempo concentram-se nas frequências mais altas do
processador, enquanto as menores energias tendem a surgir em frequências
intermediárias. Além disso, \textit{heaps} maiores que o mínimo necessário,
especialmente a partir de \(4\times\), aparecem com frequência
entre as configurações mais competitivas. Esses resultados reforçam a hipótese
inicial do trabalho e estão de acordo com a literatura que aponta que tempo e
energia não devem ser tratados como métricas intercambiáveis \cite{programming_lang_energy}, 
elas merecem ser tratadas separadamente e consideradas de forma conjunta na análise de 
desempenho.

Outro aspecto importante observado ao longo da análise é que o comportamento
energético da JVM depende da interação entre a carga de trabalho e as escolhas
de configuração. Não foi identificada uma única combinação de frequência e
\textit{heap} capaz de se destacar de forma uniforme para todos os
\textit{benchmarks}. Em vez disso, diferentes aplicações respondem de maneira
distinta às mesmas alterações, o que indica que estratégias de ajuste da JVM
baseadas em parâmetros fixos tendem a ser insuficientes quando o objetivo é
otimizar eficiência energética sem comprometer o desempenho.

Os resultados apresentados são consistentes com o objetivo proposto, porém o escopo do
estudo ainda apresenta limitações. As medições foram realizadas em um único 
\textit{hardware}, com uma única versão da JVM e um conjunto
controlado de configurações experimentais. Além disso, parte dos
\textit{benchmarks} do DaCapo Chopin não pôde ser incorporada à análise em
virtude de problemas na execução. Assim, se mostra necessário realizar essa 
análise caso a caso, considerando as características de cada aplicação e 
ambiente de execução.

Este trabalho reforça a importância de tratar a
eficiência energética de \textit{software} como um problema relevante em si,
e não apenas como reflexo de avanços em \textit{hardware}. Em contextos como
\textit{data centers}, nos quais pequenas variações de consumo podem se
acumular em custo operacional, e plataformas móveis, nas quais esse consumo
afeta diretamente a autonomia de bateria, dispor de métricas energéticas
explícitas contribui para decisões de desenvolvimento e pesquisa mais bem
fundamentadas.